% Actual class: 03
% Released material: Week3 Edit Distance
% Exact released scope: pp. 1--14
% Video supplement: Week 3 lecture transcript checked locally; not included in repository
% Human verified: Coverage and representative check answers confirmed in discussion

\section{第 03 节课：Edit Distance}
\label{lecture:SC4002-L03}

\lecturemeta
  {SC4002-L03}
  {2026-08-26}
  {Week3 Edit Distance (2026)}
  {pp.~1--14；Edit Distance 全部内容}
  {Week 3 课堂视频字幕已作本地核对，未收入仓库}
  {课堂范围及代表题答案已确认（2026-08-26）}

\subsection{本讲主线}

Edit Distance（编辑距离）用最小编辑成本衡量两个 strings（字符串）的差异，可用于 spelling correction（拼写纠错）、coreference resolution（共指消解）以及 machine translation / speech recognition evaluation（机器翻译／语音识别评估）。穷举所有 edit sequences（编辑序列）会产生大量重复子问题；Dynamic Programming（动态规划，DP）通过求解所有 prefix pairs（前缀对）得到全局最优解。

\lecturetopic{SC4002-L03-T01}{Minimum Edit Distance 与成本约定}

令 source string（源字符串）为 $X=x_1\cdots x_n$，target string（目标字符串）为 $Y=y_1\cdots y_m$。允许三种基本操作：

\begin{center}
\small
\begin{tabularx}{0.88\textwidth}{@{}>{\raggedright\arraybackslash}p{3.2cm}Xc@{}}
  \toprule
  Operation（操作） & 作用 & 本讲成本 \\
  \midrule
  Insertion（插入） & 在当前结果中插入 $y_j$ & $1$ \\
  Deletion（删除） & 删除 source 中的 $x_i$ & $1$ \\
  Substitution（替换） & 将 $x_i$ 替换为 $y_j$；相同字符视为 match（匹配） & $2$；match 为 $0$ \\
  \bottomrule
\end{tabularx}
\end{center}

Minimum Edit Distance（最小编辑距离，MED）是把 $X$ 变为 $Y$ 的所有合法 edit sequences 中 total cost（总成本）最小者。必须先确认 cost convention（成本约定）：有些资料把 substitution cost 设为 $1$，而本讲采用 $2$，考试计算应以题目或讲义给定的成本为准。

\textbf{对应例题：}\relatedexercise{SC4002-L03-E01}{\texttt{intention} 到 \texttt{execution} 的编辑成本}。

\lecturetopic{SC4002-L03-T02}{Dynamic Programming State 与 Recurrence}

定义
\[
  D[i,j]
  =\text{把 }X[1{:}i]\text{ 转换为 }Y[1{:}j]\text{ 的 minimum cost},
  \qquad 0\le i\le n,\ 0\le j\le m.
\]
所求答案为 $D[n,m]$。单位 insertion/deletion cost 下，boundary conditions（边界条件）为
\[
  D[i,0]=i,\qquad D[0,j]=j.
\]
令
\[
  \operatorname{sub}(x_i,y_j)=
  \begin{cases}
    0, & x_i=y_j,\\
    2, & x_i\ne y_j,
  \end{cases}
\]
则 recurrence（递推式）为
\[
  D[i,j]=\min\!\begin{cases}
    D[i-1,j]+1, & \text{delete }x_i,\\
    D[i,j-1]+1, & \text{insert }y_j,\\
    D[i-1,j-1]+\operatorname{sub}(x_i,y_j),
      & \text{match/substitute}.
  \end{cases}
\]

\begin{figure}[htbp]
  \centering
  \resizebox{0.86\textwidth}{!}{\begin{tikzpicture}[
  cell/.style={draw=noteBlue, rounded corners=2pt, minimum width=2.8cm,
    minimum height=1.0cm, align=center, fill=noteBlue!5},
  op/.style={-{Latex[length=2.6mm]}, thick, noteBlue},
  every node/.style={font=\small}
]
  \node[cell] (diag) at (0,1.9) {$D[i-1,j-1]$};
  \node[cell] (top)  at (4.6,1.9) {$D[i-1,j]$};
  \node[cell] (left) at (0,0) {$D[i,j-1]$};
  \node[cell, fill=ntuRed!7, draw=ntuRed] (cur) at (4.6,0) {$D[i,j]$};

  \draw[op] (diag) -- node[above, sloped, align=center]
    {match / substitute\\$x_i\rightarrow y_j$} (cur);
  \draw[op] (top) -- node[right, align=left]
    {delete\\$x_i$} (cur);
  \draw[op] (left) -- node[below]
    {insert $y_j$} (cur);

  \node[anchor=east, text=noteBlue] at (-1.65,0.95)
    {source prefix（源前缀）向下};
  \node[anchor=south, text=noteBlue] at (2.3,2.7)
    {target prefix（目标前缀）向右};
\end{tikzpicture}
}
  \caption{Edit-distance matrix（编辑距离矩阵）的三个 predecessor states（前驱状态）。行表示 source prefix，列表示 target prefix。}
  \label{fig:sc4002-l03-predecessors}
\end{figure}

递推式来自对 optimal solution（最优解）最后一步的完整分类：最后一步只能是 deletion、insertion 或 match/substitution；若此前的 prefix solution 不是最优，就可用更优前缀替换它，从而与整体最优矛盾。这体现了 optimal substructure（最优子结构）。

\lecturetopic{SC4002-L03-T03}{Backpointers、Alignment 与并列最优}

计算每个 $D[i,j]$ 时，可保存达到最小值的 predecessor（前驱）作为 backpointer（回溯指针）。从 $D[n,m]$ 回溯到 $D[0,0]$，即可恢复一个 optimal alignment（最优对齐）和对应 edit sequence。

若多个候选值相同，应保存全部并列 backpointers；一个最小距离可能对应多个不同的 optimal alignments。例如本讲成本下，一次 substitution 的成本 $2$，与一次 deletion 加一次 insertion 的成本 $1+1$ 相同，因此两类路径可能并列最优。

\textbf{对应例题：}\relatedexercise{SC4002-L03-E02}{Tie 与多个 Backpointers}。

\lecturetopic{SC4002-L03-T04}{Complexity 与空间优化}

矩阵共有 $(n+1)(m+1)$ 个 states，每个 state 只比较三个候选值，因此
\[
  \text{time complexity}=O(nm),\qquad
  \text{space complexity}=O(nm).
\]
若只需要 distance 而不恢复 alignment，每一行只依赖当前行与上一行，可用 rolling rows（滚动数组）把空间降为 $O(m)$（通常令较短字符串作为列）；若需要完整 backtrace，则仍须保存 backpointers 或采用额外的重算策略。

\subsection{一分钟回顾}

Edit Distance 把字符串转换问题写成 prefix-pair DP。$D[i,j]$ 的三个来源分别对应 deletion、insertion 和 match/substitution；边界行列表示把非空前缀与空串相互转换。回溯指针恢复具体 alignment，并列最小值意味着可能存在多个最优编辑序列。最后务必先确认 substitution cost，不能把“操作次数”和“加权总成本”混为一谈。
